\begin{problem}[Legacy problem: principal opens form a basis]
Let $A$ be a ring.  Show that the sets
\[
D(f)=\{\mathfrak p\in\Spec A:f\notin\mathfrak p\},
\qquad f\in A,
\]
form a basis for the topology on $\Spec A$.  Show also that
\[
D(f)=\varnothing
\]
if and only if $f$ is nilpotent.
\end{problem}

\paragraph{Why this problem matters.}
This is the principal legacy problem for VI/08.  It contains both defining facts that make
$D(f)$ useful: every open set can be assembled from principal opens, and the empty principal
open detects exactly nilpotence.

\paragraph{Useful ideas before starting.}
Use $D(f)=\Spec A\setminus V(f)$, the identity $D(fg)=D(f)\cap D(g)$, and
\[
\sqrt{(0)}=\bigcap_{\mathfrak p\in\Spec A}\mathfrak p.
\]

\paragraph{Guided hints.}
For a general open $U=\Spec A\setminus V(I)$, translate
$\mathfrak p\notin V(I)$ into the existence of one element $f\in I$ that is absent from
$\mathfrak p$.

\begin{solution}
First $D(1)=\Spec A$, so principal opens cover the spectrum.  Further,
\[
D(f)\cap D(g)=D(fg),
\]
because a prime avoids $fg$ exactly when it avoids both $f$ and $g$.  Hence intersections
of basis candidates are again basis candidates.

Now let
\[
U=\Spec A\setminus V(I)
\]
be an arbitrary open set.  For $\mathfrak p\in\Spec A$,
\[
\mathfrak p\in U
\iff
I\not\subseteq\mathfrak p
\iff
\text{there exists }f\in I\text{ with }f\notin\mathfrak p
\iff
\mathfrak p\in D(f)\text{ for some }f\in I.
\]
Therefore
\[
U=\bigcup_{f\in I}D(f),
\]
and the $D(f)$ form a basis.

Finally,
\[
D(f)=\varnothing
\]
iff every prime ideal contains $f$, iff
\[
f\in\bigcap_{\mathfrak p\in\Spec A}\mathfrak p=\sqrt{(0)},
\]
iff $f$ is nilpotent.
\end{solution}

\paragraph{What happened mathematically.}
The topology is generated by asking where one element can be inverted.  At the opposite
extreme, an element that cannot be inverted at any prime is exactly a nilpotent element.

\paragraph{Extensions.}
Show that $D(f)=\Spec A$ iff $f$ is a unit, and compare the two extreme cases with the
localizations $A_f$.

\paragraph{Provenance.}
Legacy \emph{Theory of Algebraic Geometry 1}, Problem 1.  The basis statement and the
nilpotence criterion are preserved together as one canonical problem.  The legacy title
also mentioned polynomial-ring spectra and disconnected spectra, but those are separate
Problems 2 and 3 in the same source and have canonical homes VI/07 and VI/05 respectively.
